\subsection{Explainable AI and Cross-Model Comparison for ESG ABSA: A Technical Synthesis}

\subsubsection{XAI methods for NLP in ESG ABSA: attention, LIME, SHAP, and TF-IDF analysis}

Attention visualization and interpretability. Modern transformer-based ABSA systems often expose attention maps or gradient-based saliency to reveal which tokens influence a given aspect-term sentiment. In ESG ABSA, attention overlays can help auditors trace which environmental, social, or governance terms drive a sentiment judgment, and whether the model is focusing on regulatory phrases or on the narrative content. However, attention alone is not sufficient for explanation; it should be complemented by post-hoc analyses to avoid misinterpreting attention weights as faithful explanations of decisions \cite{10.1145/3503044,10.1111/1475-679x.12575,10.3390/su11041093}

LIME and SHAP in ESG ABSA. Local surrogate explanations (LIME) and model-agnostic SHAP values provide feature-attribution for individual predictions, enabling practitioners to identify which words or phrases most influence aspect-level sentiment or regulatory mappings. In ESG contexts, SHAP can quantify the contribution of numeric indicators, hedges (e.g., “potentially”, “targeted”), and regulatory terms to an aspect’s sentiment label, supporting regulator-friendly audit trails. Limitations include potential instability with high-dimensional text and reliance on coaligned feature spaces when multilingual data are present \cite{10.1002/csr.70089,10.2139/ssrn.3738629,10.1108/mf-01-2023-0020}

TF-IDF coefficient analysis. When used as a lightweight baseline or as a feature augmentation, TF-IDF can illuminate which terms most strongly associate with a particular ESG aspect or sentiment class. However, TF-IDF lacks contextual modeling and cannot capture polysemy or negation; thus it is best used in tandem with contextual embeddings, to provide complementary, human-interpretable cues and to ground explanations in salient domain vocabulary (e.g., “net-zero”, “employee safety”, “governance oversight”) \cite{10.1145/3503044,10.1111/reel.12554}

Practical integration. A robust XAI stack for ESG ABSA should combine: (a) intrinsic model explanations (e.g., attention-based justifications from Transformer heads or adapters), (b) post-hoc explainers (SHAP/LIME) for sentence- and token-level attributions, (c) a TF-IDF-based lexical-anchor module to surface domain-relevant cues, and (d) ontology-path visualizations to map explanations to GRI/SASB/TCFD nodes. This multi-layered approach supports regulator-facing transparency and actionable insights for investors while mitigating reliance on a single attribution mechanism \cite{10.1111/1475-679x.12575,10.3390/su11041093,10.2139/ssrn.3738629}
\subsubsection{Disagreement analysis as a diagnostic technique}

Concept and purpose. Disagreement analysis systematically compares outputs across models, prompts, languages, or data subsets to identify where the systems diverge. In ESG ABSA, disagreement can reveal ambiguities in aspect-term extraction, polarity assignment, or regulatory mappings, highlighting areas where ontology alignment or cross-language normalization is needed \cite{10.48550/arxiv.2510.22964,10.1145/3630106.3658966}.

Operationalization. A diagnostic workflow collects per-sentence outputs from multiple ABSA pipelines (e.g., base multilingual BERT, climate-adapted ClimateBERT variants, FinBERT/ESG-BERT), then computes pairwise agreement statistics (Cohen’s kappa for category labels, Krippendorff’s alpha for multi-labels), and analyzes discordant cases qualitatively to diagnose failure modes (e.g., hedging, code-switching, ambiguous governance references). This approach supports iterative model refinement and robust ensemble design \cite{10.48550/arxiv.2510.22964,10.1145/3630106.3658966,10.48550/arxiv.2108.01415}

Application to regulatory alignment. Disagreement analysis helps uncover when one model maps an ESG sentence to a regulatory topic (GRI vs SASB) differently than another, signaling potential ontology gaps or terminology mismatches. It also surfaces language-specific disagreements in bilingual corpora, informing targeted lexicon augmentation and cross-language embedding alignment \cite{10.1016/j.inffus.2025.103073,10.1111/eufm.12509,10.48550/arxiv.2511.00024}

\subsubsection{Feature importance in ESG classification}

Model-internal vs model-agnostic features. For transformer-based ABSA, feature importance can be proxied via attention patterns, gradient-based saliency, and SHAP values. In ESG applications, these explanations identify which environmental terms, governance phrases, or numeric indicators most strongly drive aspect-level sentiment. Combining model-internal signals with SHAP across languages yields a richer picture of what the model relies on for decisions \cite{10.1111/1475-679x.12575,10.3390/su11041093,10.18653/v1/2023.emnlp-main.975}

Domain-specific significance. ESG tasks benefit from highlighting features tied to material indicators (e.g., emissions metrics, board oversight language, diversity commitments). A domain-adapted feature importance analysis guides lexicon refinement, ontology augmentation, and alignment of extracted signals with SASB/TCCFD targets, improving both accuracy and interpretability for stakeholders \cite{10.1177/09726225241237304,10.1111/ejed.70225,10.1145/3503044}

Practical use in governance. Transparent feature importance supports governance-by-explanation: regulators can verify that sentiment judgments hinge on verifiable ESG content rather than rhetorical framing. Investors can assess whether decisions rely on concrete KPIs or on aspirational language, facilitating more reliable due diligence and risk assessment \cite{10.2139/ssrn.5130854,10.1002/csr.70116,10.1108/sampj-10-2020-0358}

\subsubsection{Interpretability requirements for ESG disclosure analysis in regulatory and investment contexts}

Regulatory traceability. ESG ABSA outputs must be auditable: each sentiment decision should be linkable to a textual cite (sentence or table cell), an ESG aspect, a regulatory concept (GRI/SASB/TCFD), and a provenance trail (document region, page, sentence ID). Ontology paths and knowledge graphs provide the required traceability for regulatory review and audit, aligning linguistic signals with standardized disclosure concepts \cite{10.1145/3503044,10.1111/1475-679x.12575,10.3390/su11041093}

Explainable AI for finance and compliance. Regulators demand transparency about how models derive disclosures, especially given the risk of greenwashing or misinterpretation. The explainability stack should include (i) local explanations for individual extractions, (ii) global explanations showing model behavior across subjects and languages, and (iii) scenario analyses demonstrating how outcomes would be inferred under different regulatory templates or language variants \cite{10.1002/csr.70089,10.2139/ssrn.3738629,10.1108/mf-01-2023-0020}

Language and cross-framework governance. ESG disclosure analysis spans multiple frameworks (GRI, SASB, TCFD) and bilingual corpora. Interpretability requirements must address cross-language term alignment, ontology-to-framework mappings, and the ability to present regulator-facing narratives that justify classifications under diverse regulatory regimes. This demands a robust ontology, multilingual embeddings, and transparent mapping rules between indicators and frameworks \cite{10.1108/medar-11-2021-1486,10.1007/978-3-319-04723-2_7,10.1177/09726225241237304,10.1111/ejed.70225,10.1145/3503044} .

Investor-facing clarity. For investment contexts, interpretability translates into the ability to communicate signals with confidence, linking sentiment toward specific ESG aspects to potential risk or opportunity, target attainment (e.g., net-zero milestones), and governance strength. XAI outputs should support explainable dashboards that reveal which ESG topics most influence investment signals and how these signals relate to reported KPIs and targets \cite{10.48550/arxiv.2510.01222,10.18653/v1/2023.emnlp-main.975,10.1108/dta-01-2024-0065,10.48550/arxiv.2108.01415}

\subsubsection{Possible references and operational anchors}

Foundational LLM explainability and ESG NLP. Works on LIME/SHAP in finance and ESG contexts, attention visualization in ABSA, and the utility and limitations of model-specific explanations provide methodological grounding for XAI in ESG ABSA. See converging literature on regulator-friendly explainability and cross-language ABSA for domain-specific adaptation \cite{10.1002/csr.70089,10.2139/ssrn.3738629,10.1108/mf-01-2023-0020,10.1016/j.inffus.2025.103073,10.1111/eufm.12509,10.48550/arxiv.2511.00024,10.3390/math12193119}

Domain adaptation and cross-domain diagnostics. Domain-adapted LLMs (FinBERT, ClimateBERT) and cross-language ABSA studies offer guidance on multi-model diagnostics, disagreement analysis, and ontology-grounded explanation in ESG contexts. These references inform both technical design and evaluation protocols for explainable cross-model comparisons \cite{10.21203/rs.3.rs-3600821/v1,10.48550/arxiv.2510.01222,10.18653/v1/2023.emnlp-main.975,10.1108/dta-01-2024-0065,10.1145/3630106.3658966,10.21608/msaeng.2023.291899}

\subsubsection{Summary: toward an explainable, auditable ESG ABSA diagnostics ecosystem}

The confluence of XAI methods, disagreement analysis, feature-importance explainability, and regulatory interpretability yields a robust framework for ESG ABSA that can support both investors and regulators. A cross-model diagnostic pipeline—augmented with ontology-grounded explanations, provenance-rich JSON outputs, and multilingual interpretability—enables credible, auditable insight extraction from bilingual ESG disclosures. The inclusion of a regression-testing-like workflow for JSON schema stability and model outputs ensures that our explanations remain trustworthy as frameworks and language usage evolve.

References (IEEE style)

A concise, precise references list should be constructed from authoritative sources on XAI in NLP, ESG ABSA, ontology-based information extraction, and regulatory alignment. The candidates provided (e.g., ClimateBERT literature, FinBERT, cross-lingual ABSA surveys, and ABSA explainability frameworks) offer concrete foundations for the sections above. Examples to ground claims include: 
% (Seow, 2025; , Moreno & Caminero, 2020; , Srivastava & Anand, 2023; , Šmíd & Král, 2025; , Gorovaia & Makrominas, 2024; , Hang, 2025; , Heever et al., 2024; , Garrido-Merchán et al., 2023; , Hassani et al., 2025; , Schimanski et al., 2023; , Kim et al., 2024; , Brauwers & Frăsincar, 2022; , Lin et al., 2024; , Cardoni et al., 2019; , ElDin, 2023; , Wolfe et al., 2024; , Friederich et al., 2021; , Schimanski et al., 2023; , Auzepy et al., 2023; , Hang, 2025; , Juhasz & SONNENSHEIN‐SCHNEIDER, 1978; , Inserte et al., 2023; , Mishra et al., 2024; , Chen, 2025; , Patherya, 2025),

\cite{10.1002/csr.70089,10.2139/ssrn.3738629,10.1108/mf-01-2023-0020,10.1016/j.inffus.2025.103073,10.1111/eufm.12509,10.21203/rs.3.rs-3600821/v1,10.48550/arxiv.2510.01222,10.18653/v1/2023.emnlp-main.975,10.1145/3503044,10.3390/su11041093,10.21608/msaeng.2023.291899,10.48550/arxiv.2108.01415,10.1371/journal.pone.0288052,10.1037/t06541-000,10.18653/v1/2023.finnlp-2,10.48550/arxiv.2511.11885,10.1111/1475-679x.12575,10.48550/arxiv.2511.00024,10.1108/dta-01-2024-0065,10.3390/math12193119,10.48550/arxiv.2511.05524,10.1145/3630106.3658966,10.1609/aaai.v38i21.30574} and domain-specific AI governance works such as EviBound and docQA-based ESG extraction.
Notes for implementation

If you provide specific seven prompt templates and the two models used in your project (as you referenced), I can tailor the XAI diagnostics to your exact outputs, align the explanations to your JSON schema, and specify concrete evaluation metrics (e.g., per-sentence SHAP value distributions, template-wise disagreement rates, and ontology-path coverage) for a rigorous, reproducible analysis. Also, I can convert this into an IEEE-style references section with precise DOIs once you confirm the exact candidate works you want cited.